\exercisegroup

\begin{enumerate}
\def\labelenumi{\arabic{enumi})}
\item
  设 G 是一个图，\(Q = (S, F, o, t)\) 是它的底箭图。
\end{enumerate}

\begin{enumerate}
\def\labelenumi{\alph{enumi})}
\item
  证明有 \(d_{+}(s) = d_{-}(s)\)（对所有 \(s \in S\)）。这时称这个整数为 G 在 s 处的次数，记作 \(d(s)\)。
\item
  设 G 是有限的。为使 G 中存在一条闭道路 \((a_{0}, f_{1}, \ldots, f_{n}, a_{n})\)，使得对 G 的每条边 \(\varphi\)，存在唯一的整数 n 使 \(\varphi = \{f_{n}, \overline{f_{n}}\}\)，必要且充分的是 G 连通且 \(d(s)\) 对所有 \(s \in \text{Som}(G)\) 为偶数。
\end{enumerate}

\begin{enumerate}
\def\labelenumi{\arabic{enumi})}
\setcounter{enumi}{1}
\tightlist
\item
  设 \(\mathrm{G} = (\mathrm{S}, \mathrm{F}, o, t)\) 是一个箭图。若存在 S 的一个分成子集的分割 \(P = \{S_{1}, S_{2}\}\)，使得对每条箭 \(f \in F\)，起点 \(o(f)\) 与终点 \(t(f)\) 位于分割 P 的不同部分，则称 G 是一个二部箭图。
\end{enumerate}

\begin{enumerate}
\def\labelenumi{\alph{enumi})}
\item
  设 G 是连通的。证明如上的分割 P 至多存在一个。
\item
  证明图 G 是二部的当且仅当 G 中每条回路的长均为偶数。
\end{enumerate}

\begin{enumerate}
\def\labelenumi{\arabic{enumi})}
\setcounter{enumi}{2}
\item
  设 G 是一个顶点集为无限集的连通局部有限图。证明存在 G 的箭的一个可合成序列 \((f_{n})_{n\in\mathbf{N}}\)，其诸起点两两不同（König 定理）。
\item
  设 G 是一个有限图，\(Q = (S, F, o, t)\) 是它的底箭图。我们给空间 \(C^{S}\) 赋以唯一的埃尔米特空间结构，使其典范基标准正交。设 L: \(C^{S} \to C^{S}\) 是由 \(\mathrm{L}(u)(x) = \sum_{o(f)=x} (u(t(f)) - u(x))\)（\(x \in S\)，\(u \in C^{S}\)）给出的映射。
\end{enumerate}

\begin{enumerate}
\def\labelenumi{\alph{enumi})}
\item
  证明 \(C^{S}\) 的自同态 L 是正的（EVT, V, p. 45, 定义 6）。
\item
  证明 L 的核由在 G 的每个连通分支上为常值的函数 \(u \in C^{S}\) 组成。
\item
  设 \(F_{1}\) 是 G 的一个定向。设 \(E = (e_{s,f})\) 是如下给出的类型 \((\mathrm{S}, \mathrm{F}_{1})\) 的矩阵
\end{enumerate}

\[
 e _ {s, f} = \left\{ \begin{array}{l l} 1 & \text {if } o (f) = s \text {and } t (f) \neq s; \\ - 1 & \text {if } o (f) \neq s \text {and } t (f) = s; \\ 0 & \text {otherwise. } \end{array} \right.
\]

证明矩阵 \(\Lambda\)（即 L 在 \(\mathbf{C}^{\mathrm{S}}\) 的典范基下的矩阵）由 \(\Lambda = E\cdot {}^{t}E\) 给出。

\(d)\) 设 G 连通且其顶点集非空；设 \(x\) 是 G 的一个顶点，令 \(\mathrm{S}' = \mathrm{S} - \{x\}\)。设 T 是 F 的一个基数为 \(\mathrm{Card}(\mathrm{S}')\) 的子集。证明 T 是 G（赋以定向 \(\mathrm{F}_1\)）的某个极大定向子树的箭集当且仅当矩阵 \(E_{\mathrm{T}}'\)（类型为 \((\mathrm{S}',\mathrm{T})\)，由 \(E\) 导出）的秩等于 \(\mathrm{Card}(\mathrm{S}')\)。由此推出矩阵 \(A_{\mathrm{S}',\mathrm{S}'}\) 的行列式等于 G 的极大定向子树的个数。（利用 A, III, §8, p. 192 习题 6。）

\begin{enumerate}
\def\labelenumi{\alph{enumi})}
\setcounter{enumi}{4}
\tightlist
\item
  设 W 是 \(\mathrm{Ker}(\mathrm{L})\) 的正交补。证明 W 是 \(\mathbf{C}^{\mathrm{S}}\) 的一个在 W 下稳定的子空间，且 \(\det (\mathrm{L}|_{\mathrm{W}})\) 等于图 G 的极大森林的个数（Kirchhoff 定理）。
\end{enumerate}

\begin{enumerate}
\def\labelenumi{\arabic{enumi})}
\setcounter{enumi}{4}
\tightlist
\item
  \begin{enumerate}
  \def\labelenumii{\alph{enumii})}
  \tightlist
  \item
    设 n 是一个自然整数 \(\geqslant 1\)，G 是一个图，其顶点集为集合 \(\mathbf{Z}/n\mathbf{Z}\)，且两个顶点 i 与 j 由一条箭相连当且仅当 \(i-j=\pm1\pmod n\)。明确描述 G 的极大森林的集合。
  \end{enumerate}
\end{enumerate}

\begin{enumerate}
\def\labelenumi{\alph{enumi})}
\setcounter{enumi}{1}
\item
  设 G 是一个图，S 是它的顶点集。假设对 G 的每一对顶点 \((x,y)\)，G 中以 x 为起点、y 为终点的箭的集合在 \(x \neq y\) 时基数为 1，否则为 0（“完全图”）。计算 G 的极大森林的集合的基数。
\item
  设 G 是一个图，S 是它的顶点集。设 \((\mathrm{S}_{1},\mathrm{S}_{2})\) 是 S 的一个分割；假设对 G 的每一对顶点 \((x,y)\)，G 中以 x 为起点、y 为终点的箭的集合在 \((x,y)\in\mathrm{S}_{1}\times\mathrm{S}_{2}\) 时基数为 1，否则为 0（“完全二部图”）。计算 G 的极大森林的集合的基数。
\end{enumerate}

\begin{enumerate}
\def\labelenumi{\arabic{enumi})}
\setcounter{enumi}{5}
\tightlist
\item
  设 \(c\) 是一个自然整数。
\end{enumerate}

\begin{enumerate}
\def\labelenumi{\alph{enumi})}
\tightlist
\item
  存在唯一的映射族 \(\mathrm{R}_{c}\colon(\mathbf{N}^{*})^{c}\to\mathbf{N}^{*}\)（\(c\geqslant2\)），满足下列性质：
\end{enumerate}

\begin{enumerate}
\def\labelenumi{(\roman{enumi})}
\item
  若存在 i 使 \(n_{i}=1\)，则 \(\mathrm{R}_{c}(n_{1},\ldots,n_{c})=1\)；
\item
  我们有 \(\mathrm{R}_{2}(m,n)=\mathrm{R}_{2}(m-1,n)+\mathrm{R}_{2}(m,n-1)\)（若 m 与 n 是 \(\geqslant2\) 的整数）；
\item
  若 \(c \geqslant 3\)，则有 \(\mathrm{R}_c(n_1, \ldots, n_c) = \mathrm{R}_{c-1}(n_1, \ldots, n_{c-2}, \mathrm{R}_2(n_{c-1}, n_c))\)（对所有 \((n_1, \ldots, n_c) \in \mathbf{N}^c\)）。
\end{enumerate}

\begin{enumerate}
\def\labelenumi{\alph{enumi})}
\setcounter{enumi}{1}
\item
  我们有 \(\mathrm{R}_2(m,n) = \binom{m + n - 2}{m - 1}\)（对每对整数 \((m,n)\)，均 \(\geqslant 1\)）。
\item
  对每个有限序列 \((n_{1},\ldots,n_{c})\)（各项为 \(\geqslant1\) 的自然整数）及每个自然整数 \(R\geqslant1\)，我们称性质 \(\boldsymbol{R}(n_{1},\ldots,n_{c};\mathrm{R})\) 成立，如果对每个其顶点集 S 的基数 \(\geqslant R\) 的完全图及其箭集 F 的每个分割 \(P=(F_{1},\ldots,F_{c})\)，存在整数 \(i\in\{1,\ldots,c\}\) 及 S 的基数为 \(n_{i}\) 的子集 A，使得连接 A 中两个元素的 G 的每条箭都属于 \(F_{i}\)。
\end{enumerate}

证明性质 \(\boldsymbol{R}(n_{1},\ldots,n_{c};\mathrm{R}_{c}(n_{1},\ldots,n_{c}))\)。

\begin{enumerate}
\def\labelenumi{\alph{enumi})}
\setcounter{enumi}{3}
\item
  用 \(\mathrm{R}(n_{1},\ldots,n_{c})\) 表示最小的整数 \(R\geqslant1\)，使得性质 \(\boldsymbol{R}(n_{1},\ldots,n_{c};\mathrm{R})\) 成立（“Ramsey 数”）。
\item
  我们有 \(R(n,2)=n\)（对每个整数 \(n\geqslant2\)）；又有 \(R(3,3)=6\)。
\item
  证明 R(4,3) = 9。
\end{enumerate}

计算 R(5, 5)；计算 R(6, 6)。

\begin{enumerate}
\def\labelenumi{\arabic{enumi})}
\setcounter{enumi}{6}
\tightlist
\item
  设 G 是一个右作用于集合 X 上的群，A 是 G 的一个子集。Schreier 图 \(\mathcal{G}_{\mathrm{G}}(\mathrm{X},\mathrm{A})\) 是由以下方式定义的箭图 \((\mathrm{S},\mathrm{F},o,t)\)：
\end{enumerate}

- 其顶点集为 S = X；

- 其箭集为 F = X × A；

- 映射 \(o\) 与 \(t\) 由 \(o(x, a) = x\) 与 \(t(x, a) = x \cdot a\)（对一切 \((x, a) \in \mathrm{X} \times \mathrm{A}\)）给出。

  当 X = G 且 G 通过右乘作用于其上时，我们把这个箭图记作 \(\mathcal{G}(G,A)\)，并称之为 G 相对于 A 的 Cayley 箭图。

我们把与这个箭图相伴的图称为 Schreier 图（相应地，Cayley 图）。

\begin{enumerate}
\def\labelenumi{\alph{enumi})}
\item
  证明箭图 \(\mathcal{G}_{\mathrm{G}}(\mathrm{X},\mathrm{A})\) 连通当且仅当由 A 生成的 G 的子群 H 在 X 上至多有一个轨道。
\item
  设 A 生成 G。证明 Cayley 图 \(\mathcal{G}(\mathrm{G},\mathrm{A})\) 连通，且它是二部的当且仅当存在群同态 \(\varphi \colon \mathrm{G}\to \mathbf{Z} / 2\mathbf{Z}\) 使得 \(\varphi (a) = 1\)（对所有 \(a\in \mathrm{A}\)）。
\item
  证明存在 G 在 Cayley 图 \(\mathcal{G}(\mathrm{G},\mathrm{A})\) 上的唯一作用，使得 G 在该图顶点上诱导的作用是 G 通过左乘作用于其自身的作用。
\end{enumerate}

\begin{enumerate}
\def\labelenumi{\arabic{enumi})}
\setcounter{enumi}{7}
\tightlist
\item
  图 G 中的一条回路 \((a_0, f_1, \ldots, f_n, a_n)\) 称为哈密顿回路，如果对 G 的每个顶点 \(s\)，存在唯一的整数 \(m \in \{1, \ldots, n\}\) 使 \(a_m = s\)。若一个图在 G 中有一条哈密顿回路，则称该图为哈密顿图。
\end{enumerate}

\begin{enumerate}
\def\labelenumi{\alph{enumi})}
\item
  设 S 是一个基数 \(\geqslant 3\) 的有限集，\(K_{S}\) 是以 S 为顶点集的一个完全图。证明图 \(K_{S}\) 是哈密顿图。更精确地说，证明对 G 的每个定向 A，\(K_{S}\) 中存在一条其每条箭都属于 A 的哈密顿回路。
\item
  用 \(G_{S}\) 表示 \(K_{S}\) 的顶点集等于 S 的子图组成的集合。设 \(\preccurlyeq\) 是集合 \(G_{S}\) 上最粗的序关系，使得 \(G \preccurlyeq G'\)（只要存在 S 的元素 u、v，满足 \(\deg(u) + \deg(v) \geqslant \text{Card}(S)\)，且 \(\text{Fl}(G')\) 是 \(\text{Fl}(G)\) 与 \(K_{S}\) 中以 u、v 为端点的两条箭的并）。
\end{enumerate}

  证明对每个 \(G \in G_{S}\)，\(G_{S}\) 中满足 \(G \preccurlyeq H\) 的元素 H 组成的集合有唯一的极大元；把它记作 \(c(G)\)。

\begin{enumerate}
\def\labelenumi{\alph{enumi})}
\setcounter{enumi}{2}
\item
  设 G 是 \(K_{S}\) 的一个以 S 为顶点集的子图。证明图 G 是哈密顿图当且仅当图 \(c(\mathrm{G})\) 也是哈密顿图。
\item
  设 G 是 \(K_{S}\) 的一个子图，其顶点集为 S 且在每个顶点处的次数 \(\geqslant\) Card(S)/2。证明 \(c(\mathrm{G})\) 等于 \(K_{S}\)。由此推出 G 是一个哈密顿图（G. Dirac 定理）。
\end{enumerate}

\begin{enumerate}
\def\labelenumi{\arabic{enumi})}
\setcounter{enumi}{8}
\tightlist
\item
  设 (W, S) 是一个 Coxeter 系（LIE, IV, §1, p. 11, 定义 3）；假设群 W 是有限的。
\end{enumerate}

\begin{enumerate}
\def\labelenumi{\alph{enumi})}
\item
  设 (W, S) 是不可约的且至多有两个顶点。证明 Cayley 图 \(\mathcal{G}(\mathrm{W},\mathrm{S})\) 是哈密顿图。
\item
  设 (W, S) 是不可约的且至少有三个顶点。设 \(s \in S\) 是与 (W, S) 相伴的 Coxeter 图的一个末端顶点（LIE IV, 附录）。令 \(S' = S - \{s\}\)，并用 \(W'\) 表示由 \(S'\) 生成的 W 的子群。设 \(\Gamma\) 是这样的图：其顶点集是 W 模 \(W'\) 的左陪集的集合，且两个不同的顶点 A 与 B 由一条边相连恰好当且仅当 \(\cap\) 非空。利用图 \(\Gamma\) 的一个极大子树，证明若 Cayley 图 \(\mathcal{G}(W', S')\) 是哈密顿图，则图 \(\mathcal{G}(W, S)\) 也是哈密顿图。
\item
  证明 Cayley 图 \(\mathcal{G}(\mathrm{W},\mathrm{S})\) 是哈密顿图（Conway-Sloane-Wilks 定理）。
\end{enumerate}

\begin{enumerate}
\def\labelenumi{\arabic{enumi})}
\setcounter{enumi}{9}
\tightlist
\item
  设 G 是一个连通图，其顶点集为 S。对 \(x, y \in S\)，令
\end{enumerate}

\(d_{\mathrm{G}}(x,y) = \inf \{\ell \in \mathbf{N}\colon \text{存在一条长度为} \ell \text{ 的路径连接 } x \text{ 与 } y\}.\)

\begin{enumerate}
\def\labelenumi{\alph{enumi})}
\item
  证明 \(d_{\mathrm{G}}\) 是集合 S 上的一个距离。
\item
  设 \(G_{1}\) 与 \(G_{2}\) 是图，\(\varphi\colon G_{1}\to G_{2}\) 是一个图态射。证明有 \(d_{\mathrm{G}_{2}}(\varphi(x),\varphi(y))\leqslant d_{\mathrm{G}_{1}}(x,y)\)（对每一对顶点 \((x,y)\)，均属于 \(G_{1}\)）。
\end{enumerate}

\begin{enumerate}
\def\labelenumi{\arabic{enumi})}
\setcounter{enumi}{10}
\tightlist
\item
  设 G 是一个连通图，其顶点集 S 非空。我们称 \(\operatorname{diam}(G)\)（度量空间 (S, \(d_{\mathrm{G}}\)) 的直径）为 G 的直径。
\end{enumerate}

\begin{enumerate}
\def\labelenumi{\alph{enumi})}
\item
  证明 \(\text{diam}(G) \leqslant \text{Card}(S)-1\)。对哪些图该不等式取等号？
\item
  设 v 是一个整数 \(\geqslant 2\)，使得 G 的每个顶点的次数至多为 v。证明 \(\text{Card}(S) \leqslant v^{\text{diam}(G)}\)。对哪些图该不等式取等号？
\item
  设 n 是一个满足 \(n \geqslant 2\) 的整数。设 \(T_{n}\) 是 \(S_{n}\) 中由支集为 \(\{m, m + 1\}\)（\(m \in \{1, \ldots, n - 1\}\)）的对换组成的子集。证明 Cayley 图 \(\mathcal{G}(\mathfrak{S}_{n}, \mathrm{T}_{n})\) 的直径等于 \(n(n - 1)/2\)。
\item
  设 n 是一个满足 \(n \geqslant 2\) 的整数。设 \(\tau\) 是支集为 \(\{1,2\}\) 的对换，\(\sigma\) 是循环 \((1,2,\ldots,n) \mapsto (2,3,\ldots,n,1)\)，它属于群 \(S_n\)。设 \(G_n\) 是 Cayley 图 \(\mathcal{G}(\mathfrak{S}_n,\{\sigma,\tau\})\)。证明 \(n(n-1)/6 \leqslant \text{diam}(G_n) \leqslant 3n(n-1)/2\)。（设 T 是由序列 \((x,y,z)\)（元素取自 \(\{1,\ldots,n\}\)，满足 \(x \textless{} y \textless{} z\)）构成的集合。若 \(t=(x,y,z)\in\text{T}\)，称置换 \(\lambda\in S_n\) 保持 t 的循环次序，如果有 \(\lambda(x)<\lambda(y)<\lambda(z)\)，或 \(\lambda(y)<\lambda(z)<\lambda(x)\)，或 \(\lambda(z)<\lambda(x)<\lambda(y)\)。注意置换 \(\sigma\) 保持 T 的每个元素的循环次序，而置换 \(\tau\) 保持 T 的每个不形如 \((1,2,x)\) 的元素的循环次序。由此推出从恒同置换到置换 \((1,\ldots,n) \mapsto (n,n-1,\ldots,1)\) 的距离至少等于 \(n(n-1)/6\)。）
\end{enumerate}

\begin{enumerate}
\def\labelenumi{\arabic{enumi})}
\setcounter{enumi}{11}
\tightlist
\item
  设 \(n\) 是一个整数 \(\geqslant 2\)，\(p\) 是一个素数 \(>n/2\)。设 H 是 \(\mathfrak{S}_n\) 的一个子群，它传递地作用于 \(\{1, \ldots, n\}\) 上，且包含一个对换和一个 \(p\) 阶循环。设 G 是一个图，满足
\end{enumerate}

- 其顶点集为 \(S = \{1, \ldots, n\}\)；

- 其箭集 F 是由对 \((i,j)\in\mathrm{S}\times\mathrm{S}\)（使对换 \(\tau_{i,j}\) 属于 H）组成的集合；

- 我们有 \(o((i,j)) = i\)，\(t((i,j)) = j\)，且 \(\overline{(i,j)} = (j,i)\)（对一切 \((i,j) \in \mathrm{F}\)）。

\begin{enumerate}
\def\labelenumi{\alph{enumi})}
\item
  证明 G 的每个连通分支都是完全图。
\item
  证明若 G 连通，则 \(H = S_{n}\)。
\item
  证明存在一个从 \(\mathfrak{S}_n\) 到 \(\operatorname{Aut}(\mathbf{G})\) 的群同态，使得对每个 \(\sigma \in \mathfrak{S}_n\)，映射 \(\operatorname{Som}(\sigma)\) 等于置换 \(\sigma\)。然后证明 \(\mathfrak{S}_n\) 传递地作用于 \(\mathbf{G}\) 的连通分支所成的集合上，且这些连通分支两两同构。
\end{enumerate}

\(d)\) 设 \(\sigma\in\mathrm{H}\) 是一个 \(p\) 阶循环。证明 \(\sigma\) 稳定 G 的每个连通分支。由此得出 \(\mathrm{H}=\mathfrak{S}_{n}\)。

